\chapter*{Appendix A: QUBO Derivation Details}

\section{Maximum Weight Independent Set to QUBO}

This appendix provides the complete algebraic derivation from the MWIS problem to its QUBO representation.

\subsection{Problem Statement}

Given a graph $G = (V, E)$ with node weights $w_i > 0$ for $i \in V$, the MWIS problem is:

\begin{equation}
\max \quad \sum_{i \in V} w_i x_i
\end{equation}

subject to:
\begin{equation}
x_i + x_j \leq 1 \quad \forall (i,j) \in E
\end{equation}

where $x_i \in \{0, 1\}$.

The constraints enforce independence: if edge $(i,j)$ exists, nodes $i$ and $j$ cannot both be selected.

\subsection{Penalty Formulation}

To convert to unconstrained form, we use a penalty method. The unconstrained objective is:

\begin{equation}
\max \quad \sum_{i \in V} w_i x_i - \lambda \sum_{(i,j) \in E} x_i x_j
\end{equation}

where $\lambda > 0$ is the penalty coefficient.

Rearranging as a minimization problem:

\begin{equation}\label{eq:qubo-min-form}
\min \quad -\sum_{i \in V} w_i x_i + \lambda \sum_{(i,j) \in E} x_i x_j
\end{equation}

\subsection{Matrix Form}

Expanding the objective:

\begin{align}
f(x) &= -\sum_{i} w_i x_i + \lambda \sum_{(i,j) \in E} x_i x_j \\
&= \sum_{i} (-w_i) x_i + \sum_{(i,j) \in E} \lambda x_i x_j \\
&= \frac{1}{2} \mathbf{x}^T Q_{\text{QP}} \mathbf{x} + \mathbf{p}^T \mathbf{x}
\end{align}

where:

\begin{equation}
Q_{\text{QP},ij} = \begin{cases}
-2w_i & \text{if} \quad i = j \\
2\lambda & \text{if} \quad (i,j) \in E \\
0 & \text{otherwise}
\end{cases}
\end{equation}

and $\mathbf{p} = 0$.

---

\section{Penalty Coefficient Sufficiency Proof}

\begin{theorem}
If $\lambda > \max_{(i,j) \in E} (w_i + w_j)$, then any minimizer of Equation \eqref{eq:qubo-min-form} corresponds to an independent set.
\end{theorem}

\begin{proof}
Suppose $x_i = x_j = 1$ for some edge $(i,j) \in E$. Consider flipping one variable to zero. If we set $x_i := 0$, the change in objective is:

\begin{equation}
\Delta f = (w_i) - (0) - \lambda = w_i - \lambda
\end{equation}

Since $\lambda > w_i + w_j \geq w_i$, we have $\Delta f > 0$ (increase in objective, i.e., decrease in minimum). Thus, setting $x_i = 0$ improves the solution. At optimality, we cannot have both $x_i = 1$ and $x_j = 1$.

By contrapositive, every minimizer satisfies the independence constraint. $\square$
\end{proof}

---

\section{Worked Example Computation}

For the three-robot, two-task example:

\begin{equation}
w = [5.2094, 2.5858, 2.1487, 2.1487, 3.9692, 1.1543, 1.4633]^T
\end{equation}

The edge list is given in Chapter 4. The maximum edge weight sum is:

\begin{equation}
\max_{(i,j) \in E} (w_i + w_j) = w_0 + w_1 = 7.7952
\end{equation}

We choose $\lambda = 8.0 > 7.7952$, ensuring the penalty is sufficient.

The QUBO matrix $Q$ is $7 \times 7$:

\begin{equation}
Q_{ii} = -2w_i
\end{equation}

\begin{equation}
Q_{ij} = \begin{cases} 2\lambda = 16.0 & \text{if} \quad (i,j) \in E \\ 0 & \text{otherwise} \end{cases}
\end{equation}

The resulting QUBO problem minimizes:

\begin{equation}
f(x) = \frac{1}{2} \mathbf{x}^T Q \mathbf{x}
\end{equation}

The optimal solution to this QUBO is $x^* = [1, 0, 0, 0, 1, 0, 0]^T$, corresponding to selecting nodes $v_0$ and $v_4$, with objective value $9.1787$ (matching the analytical optimum).
